% 来源：成文/A\_10\_模型检验.md
\section{模型检验}
本章对全文模型作跨问系统检验。各问 5.x.5 只给出该问专属检验的要点，证据、图表与完整结论一律归本章，两处不复述同一结果。检验覆盖收敛性与误差带、解析对拍、独立实现互验、物理自检与断言、守恒核算、局部与全局灵敏度、预处理与简化代价、判据与经验式外推不确定度、输入不确定度与数据同化共十一类，且每一类都给出可核验的量值。为免歧义，本章所称含水率均指干基含水率。

\subsection{网格与时间收敛性（含离散误差带）}
空间取三档网格作加密序列，以广义理查森外推估计观测收敛阶与网格收敛指数；时间取三档步长作对照；对问题一的温度场另做步长无关性检验。为使结论可核，先给出两个定义式：报告值以区间给出，区间半幅取两类误差中的较大者，即

\begin{equation}
  y\in\big[\hat y-e,\ \hat y+e\big],\qquad e=\max\big(e_{\rm disc},\ e_{\rm extrap}\big)
  \label{eq:6-1}
\end{equation}
而观测收敛阶与网格收敛指数由三档序列估计：

\begin{equation}
  p=\frac{\ln\big|(f_3-f_2)/(f_2-f_1)\big|}{\ln r}
  \label{eq:6-2}
\end{equation}
三档网格与三档步长的收敛序列见图~\ref{fig:A-1}。


% 图A-1 ＝ F-06（F 号）／图05（a/b）（PDF 文件索引号）／网格与时间收敛
\begin{figure}[htbp]
  \centering
\includegraphics[width=0.485\textwidth]{../../05\_成品图/图05a-网格与时间收敛-a.pdf}  \hfill \includegraphics[width=0.485\textwidth]{../../05\_成品图/图05b-网格与时间收敛-b.pdf}
  \caption{网格与时间收敛性}\label{fig:A-1}
\end{figure}

由图~\ref{fig:A-1} 可见，空间离散按问按场呈二阶（表面含水率列为一阶），且时间侧温度场步长无关——三档序列单调、沿半径与时间未见振荡或孤立异常点，故取 0.25 mm 为交付网格。

结果：三档序列一致显示"空间二阶、时间已消除离散误差"。分场看：问题一温度场（中心／表面／均方根）观测阶 2.002／2.006／2.003，网格收敛指数不超过 0.0001\%；问题一含水率场为 2.014（均方根）与 2.028（表面列），表面列网格收敛指数 0.011\%；问题三观测阶 1.650、网格收敛指数 0.0011\%；问题四中心含水率在 24 h 上观测阶 1.98、烘干时长序列上 1.91，而表面含水率列为 0.99（一阶）。时间侧：问题三三档呈非单调，取保守上界 ±0.0802 h；问题四保守上界为 ±0.0249 h，空间项仅 0.0002 h 可以忽略；问题一温度场分段步长取 1 s 与 60 s 相差 3.1×10\textsuperscript{-11} ℃，即时间离散误差已归零，而含水率侧为后向欧拉一阶、子步 0.1 s 下误差量级仅 10\textsuperscript{-5} kg/kg。

结论：报告值已网格收敛，且其误差带是"可复算、可复现"的，而非经验估计——两类不确定度分列报告、不得相加，报告值一律采用"空间与时间误差带取较大者"的合成口径。这意味着四位小数输出的末位仍受误差支配，故正文对末位一律不作过强解读。

补充对照为三项排除性或等价性对照。其一为界面取法对照（问题一侧）：沿含水率的积分平均与沿扩散系数倒数的调和型积分给出含水率最大差 8.5×10\textsuperscript{-6} kg/kg，在四位小数下不可辨识；同一对照在问题四的干燥尾部进一步放大，这正是 5.2.3 采用积分平均的直接理由。界面取法与网格的二维对照见图~\ref{fig:A-10}。

由图~\ref{fig:A-10} 可见，错误取法下网格不收敛，这从反面确认了积分平均口径的必要性。

其二为离散分辨率扫描：步长与网格四种配置间的散布不超过 0.15\%，均非结论的主导不确定度。其三为自适应时间步：按含水率变化率在 1/64 s 至 1/4 s 之间调整内部子步，其子步数降至固定子步的 18.75\%、关键量偏差不超过 9.3×10\textsuperscript{-6} kg/kg，故报告值采用固定子步口径，以保证可复现。

\subsection{解析解对拍}
在常物性与分段常边界条件下，把问题退化为无限长圆柱的级数解问题（含真实时变边界的分段线性杜哈梅尔积分路径）；问题四另在常物性加固定半径的条件下作退化对拍。对拍目的在验证实现是否正确，而非验证模型是否完备。常物性与分段常边界下的级数解对拍见图~\ref{fig:A-2}。


\begin{figure}[htbp]
  \centering
% 图A-10 ＝ F-39（F 号）／图35（PDF 文件索引号）／界面取法与网格对照
  \begin{minipage}[t]{0.455\textwidth}
    \centering
    \includegraphics[width=1.000\linewidth]{../../05\_成品图/图35-Q3界面取法与网格对照.pdf}
    \caption{界面取法与网格对照}\label{fig:A-10}
  \end{minipage}
  \hfill
% 图A-2 ＝ F-11（F 号）／图09（PDF 文件索引号）／解析级数解对拍
  \begin{minipage}[t]{0.455\textwidth}
    \centering
    \includegraphics[width=1.000\linewidth]{../../05\_成品图/图09-解析级数解对拍.pdf}
    \caption{解析级数解对拍}\label{fig:A-2}
  \end{minipage}
\end{figure}

由图~\ref{fig:A-2} 可见，改用精确推进后偏差转为随网格二阶下降。

结果：问题一的对拍偏差为 10\textsuperscript{-3} ℃ 量级；在改用时间方向精确推进后，该偏差转为随网格二阶下降（网格步长 1／0.5／0.25 mm 下分别为 3.886／0.975／0.244×10\textsuperscript{-4} ℃），说明残差已归于空间离散；含水率侧对应为 1.029／0.267／0.083×10\textsuperscript{-3}。问题四的退化对拍最大温差为 2.4×10\textsuperscript{-10} K。

结论：实现路径与解析解一致，对拍通过；其含义是本文的离散格式与代数装配未引入系统性偏差，残差已全部归于空间离散，可继续加密网格进一步压低。适用边界亦须声明：级数解在小的傅里叶数处不适用，且仅覆盖常物性与分段常边界情形。

\subsection{独立实现互验}
另写一套独立实现（显式格式，问题三侧为 Crank–Nicolson 与 Rannacher 启动轨），在相同输入下比较全场结果；两套实现的界面系数口径已预先对齐为沿含水率的积分平均，故其差值只反映时间格式差异，而不混入口径差异。该对照以逐点全场偏差的形式完成，差异集中在时间格式、量级低于四位小数，故本项由数值承载、不再单独制图。

结果：问题一的最大温差为 1.004×10\textsuperscript{-4} ℃、含水率最大差为 1.710×10\textsuperscript{-5} kg/kg；问题三与主实现的相对差为 3×10\textsuperscript{-7}；问题四的几何退化自检与主核一致至 3.7×10\textsuperscript{-8}。

结论：两套实现给出同一答案，排除了单一代码路径的隐性缺陷与路径依赖；这意味着报告值不是某一份代码的产物，而是与实现无关的模型解。

\subsection{物理自检与看门狗}
设置方向性、单调性与上界控制三类程序化断言（共 6 项），并用人为主观注入错误的方式演示验证断言确实能捕获错误；其中极值类断言在问题四改为温度不超压缩包络上界，因压缩温升可使中心温度超过环境上限、朴素极值原理在该问不成立。问题四另设绝热不变量判决实验。

结果：全部断言通过；绝热条件下温度与半径平方之积的不变量偏差为 3.1×10\textsuperscript{-12}；人为注错均被捕获。

结论：离散格式具备可自检的数学性质——三对角矩阵严格对角占优构成 M 矩阵，使离散极值原理先有理论、后设断言，故该组断言不是事后凑数。

\subsection{守恒性核算}
对质量与能量分别做区域积分闭合核算（问题四含时变环境通量），均以相对残差表述：

\begin{equation}
  r_{m}=\frac{\big|m(t)-m(0)-\!\int_0^{t}\!\dot m_{\rm b}\,\mathrm dt\big|}{m(0)},\qquad r_{e}\ \text{同构}
  \label{eq:6-7}
\end{equation}
其理论依据为元体平衡的代数严格守恒：相邻控制体的内部界面通量在求和时两两抵消，区域总量的变化只由最外层边界通量决定，因此残差只能来自时间积分与输出截断，而不可能来自内部装配。质量与能量守恒的总账见图~\ref{fig:A-20}。


% 图A-20 ＝ F-28（F 号）／图24（PDF 文件索引号）／守恒性总账瀑布图
\begin{figure}[htbp]
  \centering
\includegraphics[width=0.660\textwidth]{../../05\_成品图/图24-Q2守恒性总账瀑布图.pdf}
  \caption{守恒性总账瀑布图}\label{fig:A-20}
\end{figure}

由图~\ref{fig:A-20} 可见，残差随问题推进未系统性增长，长时间积分的守恒性可控。

与之互为印证的是两条过程证据：温度与含水率的相轨迹轨迹形态与物性耦合一致、无越界点；全断面失水自表面向中心推进、无反向回流。

结果：问题一的质量与能量相对残差为 1.16×10\textsuperscript{-2} 与 6×10\textsuperscript{-4}（残差归因于 21 点输出加四位小数输出网格的积分截断）；问题二为 1.9×10\textsuperscript{-3} 与 1.8×10\textsuperscript{-2}；问题三为 2.4×10\textsuperscript{-3} 与 2.2×10\textsuperscript{-3}；问题四为 2.4×10\textsuperscript{-4} 与 2.0×10\textsuperscript{-4}。

结论：守恒残差均在判据之内且未系统性增长，说明长时间积分的守恒性可控；问题四的残差最小，与其守恒形式方程的形式一致性相符。该组作用于全域积分量，能暴露逐点比对看不见的实现缺陷（如能量存储项一类的装配错误），与网格收敛指数互补而非重复。

\subsection{参数灵敏度（局部）}
对关键参数作单因素扰动：扩散系数、对流换热系数、对流传质系数各扰动 ±20\%，初始含水率扰动 ±5\%，初始温度扰动 ±1 ℃；结果以龙卷风图呈现。单因素灵敏度以归一化形式报告，使不同量纲的参数可比：

\begin{equation}
  S_i=\frac{\Delta Y/Y}{\Delta X_i/X_i}
  \label{eq:6-8}
\end{equation}
其中 $Y$ 为响应量、$X_i$ 为第 $i$ 个参数。

结果：环境温度与扩散系数是烘干时程末端温度的最主要影响；对流传质系数的作用分问不同——对时程末端温度影响小，但对问题二的表面含水率仍较敏感（归一化灵敏度约 0.41 至 0.47，为该量的主控参数），而对问题三、问题四的烘干时长影响很小（约 0.15）。全部算例中结果随参数单调变化，未出现临界敏感或反号现象。

结论：结果对参数稳健，无临界敏感参数；这意味着报告值不会因参数的小幅估计误差而失去意义。按 $\pm10\%$ 与 $\pm20\%$ 两档幅度复核，上述参数影响排序与"无临界敏感参数"的判断不变，即结论不因参数估计误差的档位选择而改变。

\subsection{参数灵敏度（全局）}
采用两族独立方法（筛选法与方差分解法，后者含相关比）在参数空间内联合抽样，以检验单因素结论在参数联合变动下是否仍成立，并识别参数交互。

结果：全局方法与单因素方法给出一致的重要度排序；扩散系数的交互作用显著。问题四侧的联合抽样进一步显示，低含水率端截断因子为第一主导因子（敏感度指数 0.337，远高于扩散系数的 0.145、传质系数的 0.043、环境温度的 0.027、环境浓度的 0.021 与换热系数的 0.002；该因子与烘干时长的相关系数为 −0.833），与 6.10 的外推不确定度结论相互印证。

结论：重要度排序稳健，但参数并非彼此独立——这意味着单因素结论只能用于排序，不能直接外推为"调参收益"，本文据此不采用该外推路径。

\subsection{数据预处理对照}
对附件 1 的环境数据做平滑与不平滑两套前处理，其余设置不变。

结果：两者对温度结果的影响不超过 1.08×10\textsuperscript{-2} ℃。

结论：本文不做平滑是可接受的选择，环境数据中的实测波动不会传导为结论层面的偏差。

\subsection{简化代价对照（潜热）}
在能量方程的表面加入汽化潜热吸热项，重算并比较，以量化忽略蒸发潜热这一简化的代价。

结果：计入潜热后温度系统性地变化 0.07 至 0.10 ℃；问题四的烘干时长仅变化 0.018\%。

结论：该简化的代价已知且有界。题面未给汽化潜热与吸湿等温线，不含源项属给定信息下的必要简化而非取舍失误，其代价已在第七章列为结论的适用边界。

\subsection{判据敏感度与低含水率端外推不确定性}
对达标阈值作非对称扰动（问题三、问题四均为收紧 13.3\%，对应阈值为 0.13 kg/kg，与放宽 20\%，对应阈值为 0.18 kg/kg，两档）；对低含水率端扩散系数经验式的适用上限作多档扫描，并给出相应的概率区间。低含水率端截断因子的扫描结果见图~\ref{fig:A-9}。


% 图A-9 ＝ F-38（F 号）／图34（PDF 文件索引号）／低含水率端外推不确定度
\begin{figure}[htbp]
  \centering
\includegraphics[width=0.660\textwidth]{../../05\_成品图/图34-Q3低C端外推不确定度.pdf}
  \caption{低含水率端外推不确定度}\label{fig:A-9}
\end{figure}

由图~\ref{fig:A-9} 可见，外推段贡献 43\% 至 66\% 的时长；干燥速率随含水率的演化显示低含水率段速率骤降，解释后期自我减速；问题四侧的外推不确定度为 16.8\% 至 55.1\%，其先验属主观设定。

结果：低含水率端外推对烘干时长的贡献，问题三为 43\% 至 66\%、问题四为 16.8\% 至 55.1\%，是全文最大的一项模型不确定性；概率区间所依据的先验分布属主观设定。

结论：核心答案对判据阈值敏感——收紧 13.3\% 使烘干时长增加约 48\%、放宽 20\% 使其减少约 33\%（问题三），问题四对应为增加 28.2\% 与减少 23.4\%；该敏感性属判据选择层面，不动摇题意裁定。故本文在给出烘干时间的同时同步声明该外推区间，引用概率区间时须注明其先验为主观设定，且不得与离散误差带相加。

\subsection{输入不确定度与数据同化对照}
对附件 1 的环境时序按实测散布叠加噪声（环境温度的标准差取 0.1561 ℃），取 40 个严格样本统计结果带宽；另行数据同化（状态空间滤波与平滑，状态取水平与斜率）给出后验带宽，作为与噪声抽样相互独立的第二条路线。

结果：噪声抽样下中心温度落在 33.5393 至 33.6065 ℃、表面温度落在 36.7148 至 36.8560 ℃、表面含水率落在 1.510344 至 1.510408 kg/kg；数据同化的后验带半幅为中心温度 5.2×10\textsuperscript{-3} ℃、表面温度 1.0×10\textsuperscript{-2} ℃，方向与前者一致、带宽更窄。末位归属分场不同：温度由输入噪声支配，其量级约 10\textsuperscript{-2} ℃，约为网格误差带的 10\textsuperscript{2} 倍；而表面含水率由网格离散支配，网格误差带 ±2×10\textsuperscript{-4} 远大于噪声贡献 2.7×10\textsuperscript{-6}。

结论：两条相互独立的路线给出方向一致的带宽，说明输入不确定度的量级可信；报告值的末位不应作为精确值解读，温度末位尤其如此。
